% Part: first-order-logic
% Chapter: tableaux

\documentclass[../../../include/open-logic-chapter]{subfiles}

\begin{document}

\iftag{FOL}
      {\olchapter{fol}{tab}{ટેબ્લો}}
      {\olchapter{pl}{tab}{ટેબ્લો}}

\begin{editorial}
આ પ્રકરણ ચિહ્નિત વિશ્લેષણાત્મક ટેબ્લો પ્રણાલી રજૂ કરે છે.

ટેબ્લોને સાબિતી-પ્રણાલી તરીકે લગતી સામગ્રી સામેલ કરવા કે કાઢવા માટે
``prfTab'' ટૅગ વાપરો.
\sourcecorrection{OLTAB-001}{મૂળની \texttt{tableaux.tex}, પંક્તિ~15માં
સાબિતી-પ્રણાલી તરીકે ``natural deduction'' લખાયું છે. પ્રકરણનું શીર્ષક,
આખી સામગ્રી અને નિયામક \texttt{prfTab} ટૅગ બધા ટેબ્લો માટે છે; તેથી અહીં
``ટેબ્લો'' પુનઃસ્થાપિત કર્યું છે.}
આ આવૃત્તિમાં આ પ્રકરણનું $!A\liff!B$ એ
$(!A\lif!B)\land(!B\lif!A)$નું સંક્ષિપ્ત સ્વરૂપ છે.
આદિમ ભાષાનું કોઈ એક બંધ વિધાન~$!D$ લઈએ તો $\ltrue$ અને $\lfalse$
અનુક્રમે $!D\lor\lnot!D$ અને $!D\land\lnot!D$નાં સંક્ષિપ્ત સ્વરૂપો
છે. ટેબ્લોના નિયમો લગાડતાં પહેલાં આ સ્વરૂપોને વિસ્તારી લઈએ;
પ્રકરણના બંધ-શાખા નિયમને આ અલગ આદિમ સંકેતો માટે બદલ્યો નથી.
\end{editorial}

\olimport{rules-and-proofs}

\olimport{propositional-rules}

\iftag{FOL}{%
  \olimport{quantifier-rules}
}{}

\olimport{derivations}

\olimport{proving-things}

\iftag{FOL}{%
  \olimport{proving-things-quant}
}{}

\olimport{proof-theoretic-notions}

\olimport{provability-consistency}

\olimport{provability-propositional}

\iftag{FOL}{%
  \olimport{provability-quantifiers}
}{}

\olimport{soundness}

\iftag{FOL}{%
  \olimport{identity}
  \olimport{soundness-identity}
}{}

\OLEndChapterHook

\end{document}
